% Body fragment. Requires amsmath, amssymb, amsthm and theorem environments.
\section{A category of finite challenge contracts}
\label{sec:challenge-category}

The existing finite-demand methods can be organized without insisting that
all demands are precisions in a single linear order.  A computation may require
several coefficient identities, a denominator bound, observations at different
points, and a compatibility check.  A finite \emph{set of questions} is a more
faithful interface than their maximum.  The forward output and the backward
questions together form a category.

The relation-and-two-map pattern is related to Dialectica categories and to the
Diller--Nahm use of finite collections of challenges; de Paiva's work is an
appropriate source for that history \cite{depaiva1991}.  We use only the explicit
first-order construction proved below.  We do not import a higher-type
interpretation, a closed structure, or the totality of higher-type recursors.

\subsection{Objects, arrows, and their computational meaning}

A contract $A$ comprises two PR carrier predicates $X_A(x)$ and $U_A(u)$ and a PR
Boolean test $R_A(x,u)$, also written as the predicate $R_A(x,u)=1$.
Think of $x$ as proposed data and $u$ as a question to
which those data must give a valid answer.  Its realization assertion is
\[
 \operatorname{Real}(A)\;:\equiv\;
 \exists x\bigl(X_A(x)\land
      \forall u\,(U_A(u)\to R_A(x,u))\bigr).
\]
A contract does not assert this formula.  It specifies the formula whose
witnesses a valid construction will transport.

Use canonical codes of finite subsets of $\mathbb N$, for example strictly
increasing finite lists.  Membership, singleton, image, and finite union are PR.
The normalized representation makes equality a literal code equality; one can
instead use arbitrary lists with equality defined by membership, but must then
remember that it is a setoid equality.  Write $\operatorname{All}_A(x,S)$ for the
bounded conjunction of $R_A(x,u)$ over $u\in S$.

An arrow $A\to B$ consists of specified PR maps
\[
 f:X_A\longrightarrow X_B,\qquad
 Q_f:X_A\times U_B\longrightarrow\operatorname{Fin}(U_A),
\]
with an iPRA proof of the local law
\[
 \operatorname{All}_A(x,Q_f(x,v))\longrightarrow R_B(f(x),v)
 \quad(X_A(x),\ U_B(v)).                              \tag{C}
\]
The typing assertions say that $f(x)$ is valid and every member of the finite
query set is a valid source challenge.  All displayed type restrictions are
arithmetic guards.  The maps are total PR maps on natural codes; only their
values at valid inputs matter.

Two arrows are equal when their forward maps agree on valid witnesses and
their query sets agree on every valid pair.  This is extensional equality of the
\emph{displayed} numerical operations, expressed by first-order formulas; it is
not assumed decidable in the general presentation.  The category is initially
an external category of specified definitions and proofs, and its laws give
fixed-signature iPRA theorem schemes.  A wholly internal finite-table version
will be given below.

\begin{proposition}[Realization transport]
If $A\to B$ is a contract arrow, iPRA proves
$\operatorname{Real}(A)\to\operatorname{Real}(B)$.
\end{proposition}
\begin{proof}
Given the witness $x$ for $A$, use $f(x)$ for $B$.  For an arbitrary valid
challenge $v$ of $B$, every member of $Q_f(x,v)$ is a valid challenge of $A$,
so the hypothesis gives its $R_A$ assertion.  Finite conjunction gives
$\operatorname{All}_A(x,Q_f(x,v))$, and (C) gives $R_B(f(x),v)$.  Universal
introduction in $v$ and existential introduction with $f(x)$ finish the proof.
There is no choice of a witness separately for each $v$.
\end{proof}

\subsection{Composition and its exact demand formula}

For arrows $(f,Q_f):A\to B$ and $(g,Q_g):B\to C$, define
\[
 (g\circ f)(x)=g(f(x)),\qquad
 Q_{g\circ f}(x,w)=
       \bigcup_{v\in Q_g(f(x),w)}Q_f(x,v).             \tag{D}
\]
The union is a bounded fold over a coded finite list, followed if necessary by
sorting and removal of duplicates; hence it is PR.

\begin{theorem}[Contract composition]
Formula (D) defines a contract arrow.  With identities
$1_A(x)=x$ and $Q_{1_A}(x,u)=\{u\}$, the contracts and these arrows satisfy the
category laws up to the stated extensional equality.
\end{theorem}
\begin{proof}
Assume $R_A(x,u)$ for every $u$ in the union (D).  Fix
$v\in Q_g(f(x),w)$.  Every $u\in Q_f(x,v)$ belongs to that union, so the local
law for $f$ gives $R_B(f(x),v)$.  This holds for every such $v$, so the local law
for $g$ gives $R_C(g(f(x)),w)$.  The types of all outputs and challenges follow
from the typing clauses of the two arrows.  The identity law (C) is simply
$R_A(x,u)\to R_A(x,u)$.

Forward identity and associativity are numerical equalities.  A singleton
union proves the two query identity laws.  For three arrows $f,g,h$, the two
query sets of the associativity law both have the membership condition
\[
 u\in Q
 \quad\Longleftrightarrow\quad
 \exists w\in Q_h(g(f(x)),z)\;
 \exists v\in Q_g(f(x),w)\;
 u\in Q_f(x,v).
\]
All existential quantifiers here are bounded by displayed finite lists.
Finite-union associativity proves equality of the normalized sets.  The proof
uses induction only for finite list operations; after Booleanization their
invariants are quantifier-free.  It does not use an unbounded union, choice,
or a total evaluator for arbitrary program codes.
\end{proof}

Thus demands compose contravariantly while proposed witnesses compose
covariantly.  The formula keeps separate questions separate; when source
questions have an order and a proved domination law, a later optimization may
replace a query set by a dominating question.  Taking the maximum without that
law is not justified.

\subsection{Finite products express conjunction of tasks}

Let $A=(X,U,R)$ and $B=(Y,V,S)$.  Define their product contract by witnesses
$X\times Y$, challenges $U+V$, and
\[
 T((x,y),\operatorname{inl}u)=R(x,u),\qquad
 T((x,y),\operatorname{inr}v)=S(y,v).
\]
Pairs and tags use standard PR coding, with the appropriate guards.
The first projection sends $(x,y)$ to $x$ and sends the target question $u$ to
the singleton $\{\operatorname{inl}u\}$; the second uses the right tag.

Given arrows $(f,Q_f):C\to A$ and $(g,Q_g):C\to B$, define their pairing by
\[
 h(z)=(f(z),g(z)),\quad
 Q_h(z,\operatorname{inl}u)=Q_f(z,u),\quad
 Q_h(z,\operatorname{inr}v)=Q_g(z,v).
\]
Each local validity law is one of the laws already supplied for $f$ or $g$.
The projection equations follow from singleton union in (D).  If another arrow
has both projection equations, its forward coordinates agree with $f$ and
$g$, and its query set at every left-tagged and right-tagged question agrees
with $Q_h$.  Those are all product questions.  Hence the arrow is equal to
$h$.  This proves the categorical product universal property, not merely a
logical implication between realization assertions.

The terminal contract has one witness and no challenges.  A map to it is
unique up to our equality: its forward value is the unique witness and there
are no valid query arguments.  Consequently this category has finite
products.  No function object or exponential was needed to construct them.

\subsection{A uniform finite-table category of contracts}

There is a useful internal specialization.  An object is a pair of finite
cardinals $x,u$ and a Boolean $x\times u$ matrix.  Its witness and challenge
sets are $[x]$ and $[u]$.  An arrow to $(y,v,S)$ is a pair of finite tables
\[
 f:[x]\to[y],\qquad Q:[x]\times[v]\to[2^u],
\]
where a number below $2^u$ is the bitset of source questions.  Validity is the
bounded test
\[
 \bigwedge_{a<x}\bigwedge_{b<v}
 \left(\left[\bigwedge_{c<u}
      (c\in Q(a,b)\Rightarrow R(a,c))\right]
      \Rightarrow S(f(a),b)\right).
\]
Empty witness and challenge sets are allowed, using the finite table
conventions made explicit in the next section. For any fixed endpoints, the possible arrow
tables are finite in number, bounded by
\[
 y^x(2^u)^{xv}.
\]
Their valid sublist is obtained by bounded filtering.  Identity and
composition use table evaluation and bitset union, so all their data and laws
are uniform PR functions and iPRA theorems about natural codes.  In this case
there is an internal, decidably presented, locally finite category of all
finite contract matrices.  Its infinitely many object codes do not cause a
universal-program-evaluation problem: each evaluated arrow is a finite table
whose domain size is explicit.

\subsection{Example: addition of regular rational names}

Fix PR rational-valued families $a(x,n)$ and $b(y,n)$; each denotes one fixed
program with parameters $x$ or $y$.  Give the first the challenge relation
\[
 R_A(x,(i,j)):\quad |a(x,i)-a(x,j)|\le 2^{-i}+2^{-j},
\]
and analogously define $R_B$.  These are decidable tests on rational codes.
The product contract requires regularity of both names.  Define the output
family by the fixed PR formula
\[
 c((x,y),n)=a(x,n+1)+b(y,n+1).
\]
The forward map keeps the pair of parameters.  For target question $(i,j)$,
ask the two source questions
\[
 \{\operatorname{inl}(i+1,j+1),\;
   \operatorname{inr}(i+1,j+1)\}.
\]
Their answers imply, by the rational triangle inequality,
\begin{align*}
 |c((x,y),i)-c((x,y),j)|
 &\le |a(x,i+1)-a(x,j+1)|+|b(y,i+1)-b(y,j+1)|\\
 &\le 2(2^{-(i+1)}+2^{-(j+1)})=2^{-i}+2^{-j}.
\end{align*}
This is exactly the local arrow law.  The realization-transport theorem
therefore carries two valid regular names to a valid regular output family.
Here the displayed shifted family defines the addition operation on names;
compatibility with a chosen equivalence relation on names is a separate
representation obligation, not a conclusion of the regularity test alone.
The algebra is rational arithmetic and the demand shift is explicit.  We have
not asserted that iPRA contains an arbitrary real number as a completed set,
or a uniform evaluator for codes of all real names.  The same contract shape
applies to error bounds, derivative observations, finite covers, and algebraic
certificates whenever their tests and numerical transformations have the
required PR definitions and proved local laws.

\subsection{Variable-length chains and the induction obligation}

A variable-length chain is legitimate when there is \emph{one uniform PR}
family of step maps $f(i,x)$, query maps $Q(i,x,u)$, and contract tests
$R(i,x,u)$.  It is not legitimate merely because each entry of an unbounded
list names some separately PR program.  In the uniform case compute
\[
 x_0=x,\qquad x_{i+1}=f(i,x_i),\qquad
 S_n=\{u\},\qquad
 S_i=\bigcup_{v\in S_{i+1}}Q(i,x_i,v)\quad(i<n).
\]
The finite state list is produced by primitive recursion.  The backward
finite-set list is produced by a loop of exactly $n$ steps using reverse
indices and bounded array lookup.  This is a PR construction even if its
sets become very large.

Assume each local step law.  If $R(0,x,v)$ holds for every $v\in S_0$, then
inductively $R(i,x_i,v)$ holds for every $v\in S_i$, because each required
source query for $v\in S_{i+1}$ is a member of $S_i$.  The induction predicate
is the PR Boolean finite conjunction
\[
 \operatorname{All}(i,x_i,S_i)=1.
\]
At $i=n$ it gives $R(n,x_n,u)$.  This identifies the exact induction formula
rather than applying induction to an assertion quantifying over all
challenges.  It generalizes a scalar backward-precision schedule to any finite
collection of heterogeneous observations.  Arithmetic bounds or resource
certificates needed to define the individual step families remain part of
their local obligations.

More explicitly, validity of each $x_i$ follows from a separate QF induction
using the forward typing law; the typing law for queries then proves validity
of all members of each backward set. Apply the guarded loop rule to
$i\le n\to\operatorname{All}(i,x_i,S_i)=1$, with the fixed terminal $n$
as a parameter. Thus the recurrence is used only at indices for which the
stored state and demand arrays have been constructed.
